\banner{Design project 13: a solid-set sprinkler block for alfalfa}\label{sec:proj13}
\textbf{Brief.} A 12-ha block of alfalfa on a sandy loam (basic intake 15 mm\,h$^{-1}$) is to be irrigated by a solid-set system with sprinklers on a $12\times18$ m rectangle. The peak $ET_c$ is 8 mm\,d$^{-1}$ and the application efficiency is 0.80, so the gross depth for a 5-day interval is $8\times5/0.8=50$ mm. The designer compares four nozzles for the application rate.

\begin{table}[!htbp]\centering\small\caption{Sprinkler nozzle alternatives on a $12\times18$ m rectangular spacing.}
\begin{tabular}{rrrrrr}\toprule
\hd{Nozzle (mm)} & \hd{Pressure (kPa)} & \hd{$q$ (m$^3$/h)} & \hd{$R$ (m)} & \hd{App. rate (mm/h)} & \hd{$S_l/2R$}\\\midrule
4.0 & 300 & 1.05 & 11.8 & 4.9 & 0.51\\
4.8 & 350 & 1.64 & 13.9 & 7.6 & 0.43\\
5.6 & 350 & 2.23 & 15.1 & 10.3 & 0.40\\
6.4 & 400 & 3.11 & 17.2 & 14.4 & 0.35\\\bottomrule\end{tabular}\end{table}

\textbf{Reading.} The nozzle of 4.8 mm at 350 kPa gives 7.6 mm\,h$^{-1}$, below the intake rate of 15 mm\,h$^{-1}$, so there is no runoff; the set time for the 50 mm is $6.6$ h, and two sets a day complete the block. The number of sprinklers in operation is the block area divided by the spacing, $12\times10^4/216=555$, for a system flow of 910 m$^3$\,h$^{-1}$ if the whole block runs at once.
\fig{g13_pat}{0.7}{Radial application patterns of three sprinkler types, scaled to a wetted radius of 14 m.}
